% !TeX root = ../main.tex

\begin{itemize}
	\item Partie qui présente de manière plus globale à quoi peut servir l’intelligence artificielle dans de tels logiciels. Idée de faire de l’intelligence artificielle un assistant archiviste/documentaliste pour accélérer, dynamiser et solidifier ces logiciels face à des collections toujours plus complexes, toujours plus grandes et toujours plus hétérogènes. Revenir sur le terme de ”production” 
% \item Définition de l’intelligence artificielle 
% \item  Définition de ce à quoi on fait référence dans l’univers de gestion des collections (RAG, CV, NLP)
% \item  Faire un dialogue entre les outils d’IA existant aujourd’hui dans le monde de la gestion des collections, et les défis relevés plus tôt. 

\item \textbf{Plan :}
\subitem 1. Approfondissement du concept général d'intelligence artificielle. Quel en est le but ? Où est ce qu'on en est aujourd'hui ? 
\subitem 2. Explications des principales notions en IA : machine/deep learning notamment.
\subitem 3. Raccrochage à l'univers de la gestion des collections en me concentrant sur les domaines IA : RAG, \textit{Computer Vision} et NLP. Définition de ces domaines, de leurs objectifs, fonctionnement et de ce qu'ils regroupent notamment dans le cas du NLP qui est extrêmement vaste \textit{(Raccrochage à la notion de Information Retrieval System ou IR présenté lors de la partie 3 du chapitre précédent, NER, POS, BERT, etc)}. Article à l'appui.
\subitem 4. Reprendre les défis et difficultés cités dans la partie 2 du chapitre 1 et montrer si/comment ces solutions apportent des résolutions. 

\end{itemize}

\label{nlp}